\section{Матрица билинейной формы. Ее свойства. Преобразование матрицы билинейной формы при переходе к другому базису}

\begin{definition}
Пусть $E = \{e_1, e_2, \dots, e_n\}$ --- базис линейного пространства $V$, и $f$ --- билинейная форма на $V$. \textbf{Матрицей билинейной формы} $f$ в базисе $E$ называется квадратная матрица $A_f$ порядка $n$, элементы которой вычисляются по формуле:
\[
a_{ij} = f(e_i, e_j), \quad i, j = 1, \dots, n.
\]
В явном виде:
\[
A_f = \begin{pmatrix}
f(e_1, e_1) & f(e_1, e_2) & \dots & f(e_1, e_n) \\
f(e_2, e_1) & f(e_2, e_2) & \dots & f(e_2, e_n) \\
\vdots & \vdots & \ddots & \vdots \\
f(e_n, e_1) & f(e_n, e_2) & \dots & f(e_n, e_n)
\end{pmatrix}.
\]
\end{definition}

\begin{theorem}[О вычислении значения формы через матрицу]
Если векторы $x, y \in V$ имеют в базисе $E$ столбцы координат $X$ и $Y$ соответственно, то значение билинейной формы вычисляется по формуле:
\[
f(x, y) = X^T A_f Y.
\]
\end{theorem}

\begin{proof}
Разложим векторы по базису: $x = \sum_{i=1}^n \xi_i e_i$ и $y = \sum_{j=1}^n \eta_j e_j$. Используя билинейность $f$:
\[
f(x, y) = f\left(\sum_{i=1}^n \xi_i e_i, \sum_{j=1}^n \eta_j e_j\right) = \sum_{i=1}^n \sum_{j=1}^n \xi_i \eta_j f(e_i, e_j) = \sum_{i=1}^n \sum_{j=1}^n \xi_i a_{ij} \eta_j.
\]
В матричной форме двойная сумма записывается как произведение строки $X^T$, матрицы $A_f$ и столбца $Y$, то есть $X^T A_f Y$.
\hfill$\blacksquare$
\end{proof}

\begin{theorem}[Об изоморфизме пространства билинейных форм и пространства матриц]
Отображение, сопоставляющее каждой билинейной форме $f$ ее матрицу $A_f$ в фиксированном базисе $E$, является изоморфизмом линейного пространства всех билинейных форм на $V$ и линейного пространства квадратных матриц $M_n(F)$.
\end{theorem}

\begin{proof}
Обозначим это отображение как $\Phi: f \mapsto A_f$.
\begin{enumerate}
    \item \textbf{Линейность:} Для форм $f, g$ и числа $\lambda$ матрица формы $f+g$ имеет элементы $(f+g)(e_i, e_j) = f(e_i, e_j) + g(e_i, e_j)$, что соответствует $A_f + A_g$. Аналогично для $\lambda f$.
    \item \textbf{Инъективность:} Если $A_f = 0$, то $f(e_i, e_j) = 0$ для всех $i, j$. В силу билинейности это означает $f(x, y) = 0$ для всех $x, y$, то есть $f$ --- нулевая форма.
    \item \textbf{Сюръективность:} Для любой матрицы $A \in M_n(F)$ можно задать форму правилом $f(x, y) = X^T A Y$. Эта функция билинейна, и ее матрица в базисе $E$ в точности равна $A$.
\end{enumerate}
\hfill$\blacksquare$
\end{proof}

\begin{theorem}[О преобразовании матрицы билинейной формы]
Пусть $E$ и $E'$ --- два базиса пространства $V$, связанные матрицей перехода $T$ ($E' = E \cdot T$). Если форма $f$ имеет в базисе $E$ матрицу $A_f$, а в базисе $E'$ --- матрицу $A_{f'}$, то они связаны соотношением:
\[
A_{f'} = T^T A_f T.
\]
\end{theorem}

\begin{proof}
Пусть $X, Y$ --- столбцы координат векторов $x, y$ в базисе $E$, а $X', Y'$ --- в базисе $E'$. Известно, что $X = T X'$ и $Y = T Y'$.
Значение формы не зависит от выбора базиса, поэтому:
\[
f(x, y) = X^T A_f Y = (T X')^T A_f (T Y') = (X')^T (T^T A_f T) Y'.
\]
С другой стороны, по определению матрицы формы в новом базисе, $f(x, y) = (X')^T A_{f'} Y'$. 
Поскольку это равенство должно выполняться для любых столбцов $X'$ и $Y'$, матрицы при них должны совпадать:
\[
A_{f'} = T^T A_f T.
\]
\hfill$\blacksquare$
\end{proof}
